\scanpage{71}
\mspage{II.19}
% archive sheet number: 70
% line 1
\deleted{On a alors la suite exacte bien connue :}

% lines 2–7
2.2. On voit de $\unsure{=|\text{même}}$ que la conclusion ne change pas \deleted{$\unsure{\pi_1(X}$} si on remplace $\unsure{k|h}$ par une extension finie $k'$, $X$ par $X'=X\otimes_k k'$, et $\unsure{\xi|\eta}$ par un pt géom $\xi'$ au-dessus de $\xi$. En effet, cela revient à \deleted{\unsure{en}} remplacer $\pi_1(X,\xi)$ par un ss-groupe ouvert, image inverse du ss-groupe ouvert de $\pi_1(S,\xi)$ défini par l'extension $k'/k$ de $k$.

% lines 8–12
2.3. Supposons maintenant $X$ quasi-projectif, comme il est loisible par 2.1. Il existe alors un schéma \deleted{\unsure{quasi}} projectif $X'$ sur $k$, \deleted{\unsure{un}} tel que $X$ soit isom à un ouvert \deleted{\unsure{de}} \inserted{dense} de $X'$. On peut de plus supposer $X'$ normal (quitte à le remplacer par \inserted{le normalisé de} $X'_{\mathrm{red}}$).
% lines 13–18
Utilisant \deleted{\unsure{le}} EGA IV \unsure{17.15.14.2}, on voit qu'il \deleted{\unsure{existe}} existe une extension finie radicielle $k_1$ de $k$, telle que le normalisé \inserted{$\unsure{\tilde X'_1|\hat X'_1}$} de $X'_1\deleted{\unsure{=} }=X'\otimes_k k_1$ soit géom\textsuperscript{t} normal. \deleted{\unsure{Il en est donc}} Il en est de $=$ du normalisé \inserted{$\hat X_1$} de $X_1=X\otimes_k k_1$, qui s'identifie à \deleted{\unsure{de}} un ouvert de $\tilde X'_1$. \deleted{\unsure{il est}}
% lines 19–24
Comme $X_1$ est géom\textsuperscript{t} unibranche (EGA IV 6.15.6), $\unsure{\tilde X_1|\tilde X'_1}\unsure{-|\longrightarrow}X_1$ est radiciel surjectif, donc induit un \uline{isom} sur les $\pi_1$ (SGA 1 \uuline{VIII} \unsure{\ldots}), \inserted{et de $=$ pour $\widetilde{\bar X}_1-\bar X_1$.} Donc, conjuguant cette remarque avec 2.2, on voit qu'on est ramené au cas où $X'$ est géom\textsuperscript{t} normal.
